\section{Implementación con OpenMP}

\subsection{Del diseño manual al declarativo}

En el caso de estudio 1 la paralelización con pthreads se hizo \emph{a mano}: el programa calculaba los rangos de filas de cada hilo, creaba los $T$ hilos con \texttt{pthread\_create}, y esperaba con \texttt{pthread\_join}. Con OpenMP no se escribe nada de eso: se declara qué bucle se paraleliza y el \textit{runtime} se encarga del resto.

El punto de partida fue, como pidió el profesor, la \textbf{versión secuencial inicial}. El kernel con OpenMP vive en su propio módulo (\texttt{modules/tiempo\_mult\_openmp.c}) y el programa principal (\texttt{src/matriz\_monotonic\_openmp.c}) es una copia del secuencial con dos cambios: incluir el header del kernel nuevo y llamar a \texttt{multiplicar\_y\_medir\_openmp}.

\subsection{La directiva elegida}

El pragma que se aplica al triple bucle es:

\begin{verbatim}
#pragma omp parallel for collapse(2)
for (int i = 0; i < N; i++) {
    for (int j = 0; j < N; j++) {
        for (int k = 0; k < N; k++) {
            C[(size_t)i * N + j] +=
                A[(size_t)i * N + k] *
                B[(size_t)k * N + j];
        }
    }
}
\end{verbatim}

\subsection{Justificación a partir de los ejemplos de la especificación}

Antes de escribir el código se revisaron los ejemplos del documento oficial \textit{OpenMP Application Programming Interface Examples v6.0.1} (disponible en \texttt{github.com/OpenMP/Examples}), siguiendo la indicación del profesor de no escoger las directivas al azar sino escogerlas entendiendo. Del índice, tres ejemplos son aplicables al patrón del kernel:

\begin{center}
\begin{tabular}{|l|p{6cm}|p{6cm}|}
\hline
\textbf{Ejemplo} & \textbf{Patrón} & \textbf{Veredicto} \\
\hline
\textbf{3.1} \texttt{ploop.1.c} &
\texttt{\#pragma omp parallel for} sobre un solo loop &
Solo paraleliza el loop externo $i$. Con $N=4000$ eso da 4000 bloques de trabajo, mal repartidos entre 12 núcleos. \\
\hline
\textbf{3.8} \texttt{collapse.1.c} &
\texttt{\#pragma omp for collapse(2)} sobre dos loops anidados &
\textbf{Elegido.} Une $i$ y $j$ en un único espacio de $N^2$ iteraciones (16 millones para $N=4000$) con buen reparto entre hilos. \\
\hline
\textbf{5.3.5} \texttt{task\_dep.5.c} &
\texttt{\#pragma omp task depend(in/inout)} sobre bloques &
Descartado. Asume un kernel bloqueado y el cálculo explícito de dependencias entre bloques; es más complejo de lo que el caso requiere. \\
\hline
\end{tabular}
\end{center}

\subsection{Por qué cada elemento de la directiva}

\begin{itemize}
    \item \textbf{\texttt{parallel for} (combinada)}: crea el equipo de hilos y reparte el trabajo del loop en una sola directiva. Es más limpio que separar \texttt{\#pragma omp parallel} de \texttt{\#pragma omp for}, y el ejemplo 3.1 del PDF muestra las dos formas equivalentes.
    \item \textbf{\texttt{collapse(2)}}: es la clave de la elección. \emph{Colapsar} dos loops anidados significa que el \textit{runtime} los trata como un único espacio de iteraciones de tamaño $N \times N$, y reparte ese espacio entre los hilos. Sin \texttt{collapse}, solo se paralelizaría el loop $i$, y con $N=2000$ eso daría 2000 bloques para 12 núcleos: un reparto muy desigual, con hilos excediendo en carga y otros esperando. El loop $k$ queda secuencial dentro de cada par $(i,j)$, que es lo correcto, porque la acumulación del producto punto debe permanecer en un solo hilo.
    \item \textbf{No se usa \texttt{private(k)} explícito}: en C99 las variables de control de un \texttt{for} ya son implícitamente privadas dentro de un loop construct. Declararlo explícitamente además \emph{falla la compilación}, porque $k$ está declarada dentro del \texttt{for} y por tanto no es visible en el punto donde aparece la directiva.
\end{itemize}

\subsection{Por qué no se necesitan otras cláusulas}

Este es el punto donde el conocimiento del algoritmo es lo que garantiza la corrección:

\begin{itemize}
    \item $A$, $B$ y $C$ quedan \textbf{shared} por defecto, y eso es lo correcto: $A$ y $B$ solo se leen, y cada hilo solo escribe en $C[i \cdot N + j]$ para el par $(i,j)$ concreto que le tocó en el espacio colapsado.
    \item \textbf{No hace falta \texttt{reduction(+:)}}. Una \textit{reduction} se usa cuando varios hilos acumuran sobre la \emph{misma} variable. Aquí cada $C[i][j]$ es acumulado íntegramente por un solo hilo (el que posee ese par $(i,j)$), y los hilos distintos escriben en celdas distintas. No hay conflicto.
    \item \textbf{No hace falta \texttt{critical} ni \texttt{atomic}}. No existe región crítica: no hay ninguna variable compartida que se modifique de forma no determinista.
    \item \textbf{No hace falta \texttt{barrier}} explícita. La región paralela ya tiene una barrera implícita al final, que garantiza que todos los hilos terminaron antes de que el programa principal imprima los resultados.
\end{itemize}

\subsection{Configuración del número de hilos}

El número de hilos con el que corre la región paralela se puede fijar en tres lugares, por orden de precedencia:

\begin{enumerate}
    \item La cláusula \texttt{num\_threads()} de la propia directiva.
    \item La llamada a \texttt{omp\_set\_num\_threads()}.
    \item La variable de ambiente \texttt{OMP\_NUM\_THREADS}.
\end{enumerate}

En este trabajo se usa la variable de ambiente \texttt{OMP\_NUM\_THREADS}, que es la forma más portable y la que además permite barrer el número de hilos \emph{sin recompilar} el binario. Es la misma estrategia que mostró el profesor en clase al cambiar de 8 a 4 hilos y ver cómo cambiaba el comportamiento del programa.

\subsection{El flag de compilación no es opcional}

Un detalle fundamental de OpenMP: \textbf{las directivas se ignoran silenciosamente si no se compila con \texttt{-fopenmp}}. El programa compilaría sin errores, se ejecutaría, y daría exactamente el mismo resultado secuencial. Por eso el flag es parte del diseño y no un detalle de invocación.

Además, hay que cuidar que el flag de optimización sea el \textbf{mismo que usan las demás estrategias} de la comparación. En una primera versión de este trabajo el binario de OpenMP se compiló con \texttt{-O0} mientras que el resto de estrategias usaban \texttt{-O2}; esa diferencia hacía que el speedup de OpenMP pareciera artificialmente malo (entre 0.71$\times$ y 1.48$\times$), porque la comparación mezclaba dos efectos distintos: el paralelismo y la ausencia de optimización del compilador. Con \texttt{-fopenmp -O2}, el $T=1$ de OpenMP da 76.02\,s contra los 75.65\,s del secuencial con \texttt{-O2}: una diferencia del 0.5\%, que confirma que la base de comparación quedó justa.
